\documentclass[11pt,letterpaper]{article}

\usepackage[margin=1in]{geometry}
\usepackage{amsmath,amssymb}
\usepackage{booktabs}
\usepackage{graphicx}
\usepackage{hyperref}
\usepackage{xcolor}
\usepackage{array}
\usepackage{enumitem}
\usepackage{fancyhdr}
\usepackage{titlesec}
\usepackage{parskip}
\usepackage{microtype}

\definecolor{mainblue}{RGB}{26,74,138}
\definecolor{relayorange}{RGB}{180,80,0}
\definecolor{gatewaypurple}{RGB}{130,50,170}
\definecolor{issured}{RGB}{180,40,40}
\definecolor{okgreen}{RGB}{20,120,50}

\hypersetup{
  colorlinks=true, linkcolor=mainblue, urlcolor=mainblue, citecolor=mainblue
}

\pagestyle{fancy}
\fancyhf{}
\rhead{\small Mt.\ Washington MANET Trial 1 --- 2026-05-23}
\lhead{\small Ethan Doherty, CS Directed Study, Summer 2026}
\cfoot{\thepage}

\titleformat{\section}{\large\bfseries\color{mainblue}}{}{0em}{}[\titlerule]
\titleformat{\subsection}{\normalsize\bfseries\color{mainblue}}{}{0em}{}

\begin{document}

% ─────────────────────────────────────────────────────────────────
\begin{center}
  {\LARGE\bfseries\color{mainblue} Mt.\ Washington MANET Field Trial}\\[4pt]
  {\large Trial 1 --- 2026-05-23}\\[6pt]
  {\normalsize Ethan Doherty \quad CS Graduate Directed Study, Summer 2026}\\[2pt]
  {\small Ammonoosuc Ravine Trail (ascent) $\cdot$ Jewell Trail (descent) $\cdot$
  915\,MHz LoRa $\cdot$ Meshtastic mesh}
\end{center}

\vspace{6pt}
\hrule
\vspace{12pt}

% ─────────────────────────────────────────────────────────────────
\section{Overview}

This report documents the first empirical field trial of a Meshtastic-based LoRa mesh radio
network on Mt.\ Washington, NH. The intended objective was to characterize LoRa signal
propagation along a roughly 10-hour ascent/descent route and screen the feasibility of a
low-cost fixed-relay infrastructure for autonomous wilderness emergency communications.
The collector recorded only 6\,h\,9\,m of that route, and the ambient packets lack the
controlled endpoints and direct-hop labels needed to characterize a propagation link.

A Raspberry Pi running a custom telemetry collector served as the mesh HEAD node,
logging all received RF packets (RSSI, SNR, mesh node ID, GPS coordinates, UTC
timestamp) to an append-only JSONL stream. GPS ground truth was independently provided
by a Garmin wearable GPS unit carried simultaneously.

% ─────────────────────────────────────────────────────────────────
\section{Timeline of Events}

\begin{center}
\begin{tabular}{@{}lp{10.5cm}@{}}
\toprule
\textbf{Time (EDT)} & \textbf{Event} \\
\midrule
08:06 & Departed trailhead (Base Rd, $\approx$764\,m). Garmin GPS active. Collector not yet online. \\
09:36 & Collector came online briefly, then entered reconnect loop. \\
\textcolor{issured}{09:36 -- 12:24} & \textcolor{issured}{\textbf{Collector offline} (2\,h 48\,m). Meshtastic Python library lost serial connection and looped silently. No RF packets logged. GPS track unaffected.} \\
12:24 & Collector reconnected. Active RF telemetry resumed. \\
14:23 & Summit reached --- 1917\,m. Peak node density: $\sim$22 unique nodes in rolling 30-min window. \\
14:23 -- 18:33 & Descent via Jewell Trail. Signal quality tracked in 5-min RSSI windows. \\
18:33 & Collector stopped. The hike-map input contains 686 packet events across the
afternoon active window; the separately generated mesh catalog reports 764 RF
observations from the earlier pre-hike soak. The two totals have been reconciled
(\texttt{scripts/reconcile\_trial1\_counts.py}): they are distinct time windows on the
same decoded receiver log, not one population, so neither is a single ``nodes heard'' total. \\
18:40 & Returned to trailhead. \\
\bottomrule
\end{tabular}
\end{center}

% ─────────────────────────────────────────────────────────────────
\section{Observed RF Performance}

\begin{center}
\begin{tabular}{@{}lp{9.2cm}@{}}
\toprule
\textbf{Metric} & \textbf{Value} \\
\midrule
Packet count                  & 686 hike-map events; 764 RF observations in mesh catalog \\
Active collection window      & 12:24 -- 18:33 EDT (6\,h 9\,m) \\
Unique-node count             & 50 decoded IDs in map input; 41 RF source IDs in mesh catalog \\
Nodes with GPS position       & 24 (per mesh catalog; from received \texttt{POSITION\_APP} packets) \\
Peak concurrent nodes (30-min window) & $\sim$22 (at summit) \\
Best RSSI observed            & $-38$\,dBm \\
Analyzed 5-min descent windows with $\geq$1 packet from any source & \textcolor{okgreen}{\textbf{100\%}} \\
\bottomrule
\end{tabular}
\end{center}

The count differences above are reconciled, not sampling uncertainty: they come from
two different time windows on the same decoded receiver log. The 686 records / 50 IDs
are the afternoon hike window ($\geq$16:24\,UTC, all packet types); the 764
observations / 41 IDs are the pre-hike soak snapshot (before 02:34\,UTC, rssi-bearing
packets). The reconciliation script (\texttt{scripts/reconcile\_trial1\_counts.py})
reproduces both the 686/50 numbers and the exact 41-node catalog ID set from that single
log, and finds 18 IDs in both windows, 32 hike-only, and 23 catalog-only. In every analyzed 5-minute window of the descent interval (summit at 14:23 until the
collector stopped at 18:33), the HEAD node demodulated at least one Meshtastic packet from
\emph{some} source --- any of $\sim$50 ambient mesh nodes at unknown
distances, potentially relayed. This is evidence of ambient 915\,MHz mesh activity reaching
forested sub-alpine terrain below 1200\,m; it is \emph{not} evidence of bidirectional
coverage, of reachability to any particular gateway, or of any specific link's viability.
Those claims require the controlled-link methodology planned for Trial~2
(Section~\ref{sec:forward}).

% ─────────────────────────────────────────────────────────────────
\section{Issues Encountered}

\subsection*{Issue 1 --- Collector outage (2\,h 48\,m)}
The Meshtastic Python library (\texttt{meshtastic.serial\_interface}) lost its serial connection
to the LoRa radio and entered a silent reconnect loop. No exception was raised; the gap
was discovered post-hike via timestamp analysis of the JSONL stream. This is a
software/serial-path outage in this record; the available evidence does not isolate the
library as the sole cause or rule out a hardware/USB contribution. The gap covered a
substantial part of the Ammonoosuc Ravine ascent, so that interval has GPS ground truth
but no RF observations.

\subsection*{Issue 2 --- No GPS pairing before departure}
The operator's phone GPS was not paired with the Meshtastic network before leaving the
parking lot. The HEAD node's own position was therefore not broadcast into the mesh
during the trial. The Garmin track later supplies receiver position, but it cannot supply
the original transmitter position, direct-hop identity, or synchronized link geometry.

\subsection*{Issue 3 --- No hop-count ground truth}
The collector version deployed during Trial~1 predated the addition of
\texttt{hop\_limit} and \texttt{hop\_start} fields. It is therefore not possible to
determine whether any given packet was a direct link or a mesh relay hop. Trial~2 will
address this with a dual-radio setup (see Section~\ref{sec:forward}).

% ─────────────────────────────────────────────────────────────────
\section{Propagation Model}

\subsection{Log-Distance Path Loss}

The standard log-distance path loss model used across LoRa mountain SAR research
\cite{bianco2021, review2024} is:

\begin{equation}
  PL(d) = PL(d_0) + 10n \log_{10}\!\left(\frac{d}{d_0}\right) + X_\sigma
  \label{eq:logdist}
\end{equation}

where $d_0$ is a reference distance, $n$ is the path loss exponent, and $X_\sigma$ is a
zero-mean Gaussian random variable with standard deviation $\sigma$ (dB) representing
log-normal shadowing. An alternative floating-intercept form \cite{review2024} is:

\begin{equation}
  PL(d) = \alpha + 10\beta \log_{10}(d) + X_\sigma
  \label{eq:floating}
\end{equation}

where $\alpha$ and $\beta$ are fitted from measurements. The two forms are related by
$\beta = n$ and $\alpha = PL(d_0) - 10\beta\log_{10}(d_0)$; anchoring the reference
distance in free space gives $PL(d_0) = 20\log_{10}(4\pi d_0/\lambda)$.

\medskip

Bianco et al.\ \cite{bianco2021} measured $n \approx 2.0$--$3.5$ in mountainous terrain
with $\sigma \approx 6$--$8$\,dB, and found that standard propagation models
significantly underpredict path loss at 915\,MHz relative to measurements, with $\alpha$
values deviating by up to 10\,dB over the full LoRa radio range.

\subsection{Free-Space Path Loss Baseline}

The free-space path loss ($n = 2$, $X_\sigma = 0$) at 915\,MHz serves as the lower
bound on path loss in this model:

\begin{equation}
  FSPL(d) = 32.44 + 20\log_{10}(d_{\mathrm{km}}) + 20\log_{10}(f_{\mathrm{MHz}})
  \label{eq:fspl}
\end{equation}

At 915\,MHz this evaluates to $FSPL(d) = 91.65 + 20\log_{10}(d_{\mathrm{km}})$\,dB.

\subsection{Terrain Loss Augmentation}

Because standard propagation models are insufficient for heterogeneous mountain terrain
\cite{review2024}, an empirical terrain-class loss term $L_T$ is added to
Eq.~\eqref{eq:fspl}:

\begin{equation}
  PL_{\mathrm{aug}}(d, z) = FSPL(d) + L_T(z)
\end{equation}

where the terrain class is determined by elevation $z$:

\begin{equation}
  L_T(z) =
  \begin{cases}
    3\,\mathrm{dB}  & z \geq 1500\,\mathrm{m} \quad (\text{alpine, open}) \\
    15\,\mathrm{dB} & 1200 \leq z < 1500\,\mathrm{m} \quad (\text{sub-alpine transition}) \\
    25\,\mathrm{dB} & z < 1200\,\mathrm{m} \quad (\text{forested})
  \end{cases}
  \label{eq:terrain}
\end{equation}

The forested value of 25\,dB is consistent with the body of literature on
vegetation attenuation at 900\,MHz-band frequencies, where canopy excess loss of
20--30\,dB is commonly reported for dense mixed forest at ranges beyond 1\,km
\cite{review2024}.

\subsection{Received Signal Power and Link Budget}

The predicted received signal power is:

\begin{equation}
  P_r(d, z) = P_t + G_t + G_r - PL_{\mathrm{aug}}(d, z)
  \label{eq:rxpower}
\end{equation}

where $P_t = 22$\,dBm (Heltec LoRa32 V3 maximum TX power), and $G_t = G_r = 2.15$\,dBi
(dipole antenna gains). The total link budget is:

\begin{equation}
  \mathcal{L} = P_t + G_t + G_r - P_{\mathrm{sens}} = 22 + 2(2.15) - (-131) = 157.3\,\mathrm{dB}
  \label{eq:budget}
\end{equation}

The receiver sensitivity $P_{\mathrm{sens}} \approx -131$\,dBm corresponds to SF11,
BW\,=\,250\,kHz, CR\,=\,4/5 --- the Meshtastic \textit{LongFast} modem preset used during
this trial (Meshtastic documents a 153\,dB link budget at 22\,dBm TX with 0\,dBi antennas,
implying $P_{\mathrm{sens}} \approx -131$\,dBm, consistent with the Semtech SX1262
datasheet). Note that SF12/BW\,125\,kHz is the \emph{deprecated LongSlow} preset, not
LongFast.

\subsection{LoRa Physical Layer Parameters}

The data rate of a LoRa symbol is determined by the spreading factor and bandwidth
\cite{pdr2021}:

\begin{equation}
  R_s = \frac{BW}{2^{SF}}
  \label{eq:symrate}
\end{equation}

\begin{equation}
  DR = SF \cdot \frac{BW}{2^{SF}} \cdot CR
  \label{eq:datarate}
\end{equation}

Each unit increase in SF provides approximately 2--3\,dB improvement in SNR sensitivity
at the cost of a halved data rate \cite{pdr2021}. The LongFast preset (SF11,
BW\,=\,250\,kHz, CR\,=\,4/5) yields a data rate of $\approx$1.07\,kbps and a sensitivity
floor of $\approx -131$\,dBm; the deprecated SF12/BW\,125\,kHz LongSlow preset would
yield $\approx$0.18\,kbps at $\approx -137$\,dBm.

\subsection{Packet Delivery Ratio}

Bianco et al.\ \cite{bianco2021} reported that the path loss at the 50\% packet delivery
ratio (PDR) threshold is approximately 135\,dB in mountainous canyon terrain.
The measured range at $\geq$50\% PDR was 300\,m in worst-case canyon geometry --- more
than five times the operating range of standard ARVA avalanche beacons under the same
conditions. In open alpine terrain, ranges extended to several kilometres.

The effective signal power for LoRa when $SNR < 0$\,dB is given by \cite{review2024}:

\begin{equation}
  ESP = RSS + SNR - 10\log_{10}\!\left(1 + 10^{0.1 \cdot SNR}\right)
  \label{eq:esp}
\end{equation}

% ─────────────────────────────────────────────────────────────────
\section{Historical Relay Screen (Superseded by Terrain-Profile Analysis)}

\subsection{Node Placement}

Three candidate nodes were initially selected at GPS-derived treeline locations using an
elevation-class/FSPL heuristic. The original assertion that each location had simultaneous
line-of-sight uphill and downhill was not based on a terrain-profile test and is false for
the summit-facing links; Section~\ref{sec:itm} supersedes it.

\begin{center}
\begin{tabular}{@{}lllll@{}}
\toprule
\textbf{Node} & \textbf{Role} & \textbf{Latitude} & \textbf{Longitude} & \textbf{Elev.} \\
\midrule
\textcolor{relayorange}{Ammo Relay}    & Treeline relay    & 44.26616°N & 71.32348°W & 1201\,m \\
\textcolor{relayorange}{Jewell Relay}  & Treeline relay    & 44.28376°N & 71.33583°W & 1199\,m \\
\textcolor{gatewaypurple}{Trailhead Gateway} & Starlink uplink & 44.26700°N & 71.36083°W & 764\,m \\
\bottomrule
\end{tabular}
\end{center}

\subsection{Historical FSPL-Screen Predictions}

Applying Eq.~\eqref{eq:rxpower} with the terrain augmentation of Eq.~\eqref{eq:terrain}:

\begin{center}
\begin{tabular}{@{}lllll@{}}
\toprule
\textbf{Link} & \textbf{Distance} & \textbf{Terrain} & \textbf{Predicted RSSI} & \textbf{Margin above $P_\mathrm{sens}$} \\
\midrule
Ammo Relay $\to$ Summit  & 1.68\,km & Alpine  & $-72.9$\,dBm & 58.1\,dB \\
Jewell Relay $\to$ Summit & 2.98\,km & Alpine  & $-77.9$\,dBm & 53.1\,dB \\
Ammo Relay $\to$ Gateway  & 2.98\,km & Forest  & $-99.8$\,dBm & 31.2\,dB \\
Jewell Relay $\to$ Gateway & 2.73\,km & Forest & $-99.1$\,dBm & 31.9\,dB \\
\bottomrule
\end{tabular}
\end{center}

Within this historical non-geometric screen, all four links clear the assumed receiver
sensitivity floor ($-131$\,dBm) by more than 31\,dB. The earlier analysis subtracted
$3\sigma$ using a shadowing parameter reported in another Alpine study
\cite{bianco2021}; that does not establish a 99.87\% reliability level for these sites.
The sites, measurement protocol, residual distribution, canopy, and terrain blockage
were not shown to match, and the subsequent terrain-profile results reverse the summit
conclusion.

\textbf{Model limitation.} Eq.~\eqref{eq:rxpower} with the elevation-class loss of
Eq.~\eqref{eq:terrain} contains no terrain-profile geometry: no line-of-sight test, no
Fresnel-zone clearance, and no knife-edge diffraction. These link predictions are
therefore \emph{screening estimates only}. Section~\ref{sec:itm} re-evaluates every
proposed link with the Longley--Rice ITM model over real USGS 3DEP terrain --- and
\textbf{reverses the conclusions of this table for the summit-facing links}.

\subsection{Maximum Planning Range}

Range is computed against an explicit \emph{planning threshold}
$P_{\mathrm{thr}} = -100$\,dBm rather than the absolute sensitivity floor. The
$\approx$31\,dB standoff from $P_{\mathrm{sens}} \approx -131$\,dBm provides a fade
margin covering $>3\sigma$ log-normal shadowing ($\sigma \approx 7$\,dB
\cite{bianco2021}) plus interference and hardware aging. Setting
$P_r = P_{\mathrm{thr}}$ in Eq.~\eqref{eq:rxpower}, with
$PL_{\max} = P_t + G_t + G_r - P_{\mathrm{thr}} = 126.3$\,dB:

\begin{equation}
  d_{\max} = 10^{\,\frac{PL_{\max} - L_T - 91.65}{20}}\,\mathrm{km}
  \label{eq:dmax}
\end{equation}

\begin{align*}
  d_{\max}^{\,\text{alpine}} &= 10^{(126.3 - 3 - 91.65)/20} = 38.2\,\mathrm{km} \\
  d_{\max}^{\,\text{forest}} &= 10^{(126.3 - 25 - 91.65)/20} = 3.0\,\mathrm{km}
\end{align*}

(Solving instead for $P_r = P_{\mathrm{sens}}$ gives FSPL-only ranges exceeding
1{,}000\,km --- a physically meaningless extrapolation, since free space with a flat
terrain-class constant ignores diffraction and terrain blockage entirely; this is why
a planning threshold is used.) This historical model produced the statement that the
2\,h 48\,m collector gap ``would have been covered.'' That is not a defensible coverage
claim: it ignores the blocking terrain that causes the ITM result in
Section~\ref{sec:itm}. The later model screen is itself unvalidated and does not replace
field PDR measurement.

\subsection{Automation Chain}

The proposed concept would require the following chain:
\[
  \underbrace{\text{Hiker triggers SOS}}_{\text{Meshtastic button}}
  \;\xrightarrow{\text{LoRa hop}}\;
  \underbrace{\text{Relay node}}_{\text{treeline}}
  \;\xrightarrow{\text{LoRa hop}}\;
  \underbrace{\text{Gateway Pi}}_{\text{trailhead}}
  \;\xrightarrow{\text{Starlink}}\;
  \underbrace{\text{NH Fish \& Game dispatch}}_{\text{SMS / email}}
\]

Target end-to-end latency: $<$30\,seconds. Trial~1 did not implement or test this chain,
gateway reachability, dispatch integration, acknowledgments, authorization, or agency
workflow, so the target is a requirement rather than a result. The historical rough BOM
was $\approx\$690$ (two relay nodes at $\sim\$120$ each, gateway Pi + radio $\sim\$150$,
Starlink Mini $\sim\$300$); prices, power hardware, mounting, weatherization, service,
permits, maintenance, and any existing backhaul must be quoted before this is presented
as a deployment cost.

% ─────────────────────────────────────────────────────────────────
\section{Terrain-Profile Verification (ITM/Longley--Rice)}
\label{sec:itm}

Every proposed link was re-evaluated with the Longley--Rice Irregular Terrain Model
(ITM v7.0, point-to-point mode) via the \texttt{itmlogic} reference implementation
\cite{itmlogic2020}, using terrain profiles sampled at 15\,m from a $\sim$7\,m/px USGS
3DEP elevation export. Parameters: 915\,MHz, vertical polarization, continental
temperate climate, $N_s = 301$, $\epsilon_r = 15$, $\sigma = 0.005$\,S/m; antenna
heights 3\,m (relay/gateway masts) and 1.5\,m (hiker handheld); 50\% confidence.
Predicted RSSI $= P_t + G_t + G_r - L_b(q)$ at reliability quantiles $q$.
Pipeline: \texttt{scripts/dem\_3dep.py} + \texttt{scripts/itm\_relay\_links.py};
artifacts in \texttt{artifacts/itm/}.

\subsection{Proposed links over real terrain}

\begin{center}
\small
\begin{tabular}{@{}llrrrl@{}}
\toprule
\textbf{Link} & \textbf{Dist} & \textbf{RSSI $q_{50}$} & \textbf{RSSI $q_{90}$} &
\textbf{F$_1$ clear.} & \textbf{Verdict} \\
\midrule
Ammo Relay $\to$ Summit    & 1.68\,km & $-132.2$ & $-132.3$ & $-5.34$ & \textcolor{issured}{\textbf{below sensitivity}} \\
Jewell Relay $\to$ Summit  & 2.98\,km & $-118.7$ & $-118.7$ & $-1.90$ & \textcolor{issured}{marginal (no margin)} \\
Ammo Relay $\to$ Gateway   & 2.98\,km & $-76.6$  & $-76.6$  & $+0.89$ & \textcolor{okgreen}{\textbf{strong}} \\
Jewell Relay $\to$ Gateway & 2.73\,km & $-74.1$  & $-74.1$  & $+1.19$ & \textcolor{okgreen}{\textbf{strong}} \\
Ammo $\leftrightarrow$ Jewell Relay & 2.19\,km & $-130.1$ & $-130.2$ & $-4.51$ & \textcolor{issured}{\textbf{below margin}} \\
\bottomrule
\end{tabular}
\end{center}

\noindent (F$_1$ clear.\ = worst first-Fresnel-zone clearance fraction with 4/3
effective Earth radius; $\geq 0.6$ is the usual engineering criterion, negative means
the geometric ray is below terrain.)

\textbf{ITM reverses the screening model.} The FSPL + terrain-class table in
Section~6.2 predicted the summit links as the strongest ($-72.9$\,dBm) and the
forested gateway links as the weakest ($-99.8$\,dBm). Over real terrain the opposite
holds: the valley-crossing gateway links are clear-LOS and strong, while the
treeline-to-summit chord is blocked by the \emph{convexity of the summit cone itself}
--- terrain rises up to 53\,m above the direct ray at $\sim$1.2\,km from the Ammo
relay. Mitigations tested and rejected: a 10\,m relay mast ($-131.9$\,dBm; the
obstruction is mid-path, not local), relocating to the Lakes of the Clouds shelf at
1{,}529\,m ($-121.5$\,dBm to summit, still diffracted), and a 10\,m summit-tower node
($-113.2$\,dBm to the gateway --- blocked in the other direction). At 915\,MHz with
practical antenna heights, \emph{no single-hop link from below clears the upper cone}.

\subsection{Collector-gap segment, re-examined}

Sampling the Garmin track across the 2\,h\,48\,m collector gap (09:36--12:24 EDT,
84 points at 120\,s) and taking the best of the two proposed relays per point (ITM
$q_{90}$): \textbf{100\%} of points are predicted above raw sensitivity
($-131$\,dBm), but only \textbf{45.2\%} clear the $-100$\,dBm planning threshold, and
33\% have geometric line-of-sight; the worst point is $-121$\,dBm at 3.07\,km. The
Section~6.3 claim that the gap ``would have been covered'' must therefore be
downgraded: coverage of the Ammonoosuc Ravine ascent is \emph{marginal} ---
physically plausible at the assumed sensitivity but with zero fade margin and
unmodeled canopy loss below treeline. It cannot be claimed as robust.

\textbf{This disagreement is itself a measurable hypothesis.} The screening model and
ITM differ by up to $\sim$60\,dB on the summit links. Trial~2 will place a beacon
node at the proposed Ammo relay coordinates and measure PDR while ascending the
ravine (Section~\ref{sec:forward}); the observed PDR-vs-position profile
discriminates directly between the two models and grounds any deployment proposal in
measurement rather than prediction.

\subsection{Airtime and SOS-chain capacity}

LongFast (SF11, BW\,250\,kHz, CR\,4/5, 16-symbol preamble): $T_{\mathrm{sym}} =
2^{11}/250\,\mathrm{kHz} = 8.192$\,ms. Computed airtimes (Semtech AN1200.13
formulation; \texttt{scripts/lora\_airtime.py}): position packet ($\sim$40\,B)
559\,ms, SOS text ($\sim$64\,B) 723\,ms. A 3-hop SOS chain is $\sim$2.2\,s of serial
airtime; with per-hop CSMA backoff under light load, end-to-end delivery remains
seconds --- the $<$30\,s target has order-of-magnitude margin \emph{when the channel
is lightly loaded}. The caveat is load: one position beacon per node per 60\,s costs
0.93\% airtime per node, and the $\sim$22 concurrent nodes observed at the summit
imply $\sim$20\% raw channel utilization before mesh rebroadcast amplification.
Congestion management (beacon intervals, hop limits, ideally a dedicated channel for
infrastructure traffic) must be part of any deployment plan.

% ─────────────────────────────────────────────────────────────────
\section{Plans Going Forward}
\label{sec:forward}

\subsection*{Trial 2 --- Hop-count ground truth}
Two Heltec V3 radios may be connected to the Pi simultaneously for redundancy:
\begin{itemize}[noitemsep]
  \item \textbf{Radio A}: restricted to \texttt{maxHops=1}
  \item \textbf{Radio B}: unrestricted (\texttt{maxHops=3})
\end{itemize}
The two-radio comparison does \emph{not} by itself establish hop truth: independent
packet loss can make a direct packet appear on only one receiver, and a hop-limit
setting is not a direct-link label. Calibration eligibility is determined per packet
from recorded \texttt{hop\_start}/\texttt{hop\_limit} metadata with
\texttt{hops\_away == 0}, a dedicated trial channel, monotonic transmitter sequence
numbers, and surveyed endpoints. Receiver logs must also record which physical radio
produced each observation.

\subsection*{Trial 2 --- Pre-departure checklist}
\begin{enumerate}[noitemsep]
  \item Pair phone GPS before leaving parking lot; verify \texttt{lat}/\texttt{lon} updating in JSONL
  \item Place static relay nodes at known waypoints with recorded coordinates
  \item Run \texttt{head\_readiness\_report.py} from parking lot before departure
  \item Confirm no restart timers active on the Pi
\end{enumerate}

\subsection*{Infrastructure proposal}
Engage NH Fish \& Game and Mt.\ Washington State Park to evaluate a pilot deployment
of relay nodes at the treeline crossings identified above.

\subsection*{Research outreach}
Contact authors of the relevant Alpine LoRa studies only after verifying their current
affiliations and contact details. Their measurements \cite{bianco2021, bianco2020}
provide methodological comparison data; they do not calibrate or validate this
project's terrain model.

% ─────────────────────────────────────────────────────────────────
\section{Calibration Methodology and Trial~1 Findings}

The AIRMap calibration pipeline was implemented and executed against the Trial~1
telemetry. The pipeline computes a free-space path loss baseline (Eq.~\ref{eq:fspl}),
joins predicted and observed signal metrics by timestamp, and fits the
floating-intercept model of Eq.~\ref{eq:floating} by ordinary least squares on
\emph{calibration-eligible} rows only (GPS on both ends, verified direct link,
plausible received power, fresh position fixes), using effective signal power (ESP,
Eq.~\ref{eq:esp}) as the observable where SNR is available. Evaluation uses
contiguous-time blocked cross-validation; parameter uncertainty is estimated with a
moving-block bootstrap. A constant mean-bias correction is retained only as a
comparison model. The Garmin GPS track was used as HEAD position ground truth.

\subsection*{Why Trial~1 did not produce valid calibration data}

Post-hoc analysis of the calibration dataset identified three structural measurement
gaps that rendered the output numerically unreliable:

\begin{enumerate}[noitemsep]
  \item \textbf{No hop-count field.} The collector version deployed in Trial~1 did
    not capture \texttt{hop\_limit} or \texttt{hop\_start}. Meshtastic logs the RSSI
    of the \emph{last hop}, not the original transmitter. Without hop count it is
    impossible to filter to direct links, making it impossible to associate an
    observed RSSI with a known source-to-receiver distance.

  \item \textbf{Adjacent device forwarding.} A second Heltec radio
    (\texttt{!db51af80}) was carried co-located with the HEAD node. This device
    forwarded packets from distant Meshtastic users across NH/VT/ME. The HEAD logged
    these forwarded packets with the original sender's GPS coordinates (55--183\,km
    away) but the adjacent device's RSSI ($-1$ to $-18$\,dBm). The resulting
    distance--RSSI pairs are physically invalid and contaminate the calibration
    dataset.

  \item \textbf{No independently placed static node.} The trial did not include a
    second radio at a fixed, independently measured GPS location. Without a
    controlled source-to-receiver link, there is no ground-truth distance--RSSI
    pair against which to calibrate.
\end{enumerate}

\subsection*{What Trial~2 will produce}

Trial~2 deploys \texttt{meshnode1} as a controlled transmitter at a surveyed site
while the HEAD follows a GPS-tracked route. Only observations that pass direct-hop,
endpoint-position, time-sync, sequence-denominator, and plausibility gates enter
calibration. Other packets remain in the raw dataset and are labelled ineligible.
The feasible sampling amendment and exact opportunity accounting are recorded in
\texttt{docs/trial2-preregistration.md}.

\begin{thebibliography}{9}

\bibitem{bianco2021}
G.\ M.\ Bianco, R.\ Giuliano, G.\ Marrocco, F.\ Mazzenga, and A.\ Mejia-Aguilar,
``LoRa System for Search and Rescue: Path-Loss Models and Procedures in Mountain
Scenarios,''
\textit{IEEE Internet of Things Journal}, vol.\ 8, no.\ 3, pp.\ 1985--1999, 2021.
DOI: 10.1109/JIOT.2020.3017044

\bibitem{bianco2020}
G.\ M.\ Bianco, A.\ Mejia-Aguilar, and G.\ Marrocco,
``Performance Evaluation of LoRa LPWAN Technology for Mountain Search and Rescue,''
in \textit{Proc. 5th International Conference on Smart and Sustainable Technologies
(SpliTech)}, 2020, pp.\ 1--4.
DOI: 10.23919/SpliTech49282.2020.9243817

\bibitem{review2024}
J.~A.\ Azevedo and F.\ Mendonça,
``A Critical Review of the Propagation Models Employed in LoRa Systems,''
\textit{Sensors}, vol.~24, no.~12, p.~3877, 2024.
DOI: 10.3390/s24123877

\bibitem{pdr2021}
Q.\ Guo, F.\ Yang, and J.\ Wei,
``Experimental Evaluation of the Packet Reception Performance of LoRa,''
\textit{Sensors}, vol.~21, no.~4, p.~1071, 2021.
DOI: 10.3390/s21041071

\bibitem{itmlogic2020}
E.~J.\ Oughton, T.~Russell, J.~Johnson, C.~Yardim, and J.~Kusuma,
``itmlogic: The Irregular Terrain Model by Longley and Rice,''
\textit{Journal of Open Source Software}, vol.~5, no.~51, p.~2266, 2020.
DOI: 10.21105/joss.02266

\end{thebibliography}

\end{document}
